\lecture{10}{Градиентный спуск: двойственность}

\sourcevideo
    {https://www.youtube.com/playlist?list=PLOzRYVm0a65fhKDS8cYIC7YCuVGJ4FOJx}
    {Week 3: Lecture 10: Analysis of Gradient Descent Algorithm}

В начале лекции обсуждается влияние шага градиентного спуска на
сходимость. Основная часть посвящена построению двойственных задач
для линейного программирования и метода опорных векторов. Двойственная
форма SVM приводит к ядровому методу, позволяющему строить нелинейные
границы классификации.

\section{Гладкость, сильная выпуклость и градиентный спуск}

Пусть \(f\colon\RR^n\to\RR\) дифференцируема.

\begin{definition}[\(L\)-гладкость]
Функция \(f\) называется \(L\)-гладкой, если её градиент
\(L\)-липшицев:
\[
    \norm{\nabla f(x)-\nabla f(y)}
    \le
    L\norm{x-y}
\]
для любых \(x,y\in\RR^n\).
\end{definition}

Если \(f\in C^2(\RR^n)\), то \(L\)-гладкость эквивалентна оценке
\[
    \lVert\nabla^2 f(x)\rVert_{\mathrm{op}}
    \le L
    \qquad
    \text{для всех }x\in\RR^n.
\]
В частности, для выпуклой функции это условие имеет вид
\[
    0\preceq\nabla^2 f(x)\preceq LI.
\]

Для \(\mu\)-сильно выпуклой дважды дифференцируемой функции
\[
    \nabla^2f(x)\succeq\mu I.
\]
Следовательно, если функция одновременно \(L\)-гладка и
\(\mu\)-сильно выпукла, то
\[
    \mu I
    \preceq
    \nabla^2f(x)
    \preceq
    LI,
\]
откуда
\[
    0<\mu\le L.
\]

Отношение \(\mu/L\in(0,1]\) характеризует относительную нижнюю
границу кривизны. Величину
\[
    \kappa=\frac{L}{\mu}\ge1
\]
называют числом обусловленности функции на рассматриваемом классе:
чем больше \(\kappa\), тем сильнее различаются масштабы кривизны и тем
медленнее обычно сходятся методы первого порядка.



Для задачи
\[
    \minimize_{x\in\RR^n}f(x)
\]
градиентный спуск задаётся правилом
\[
    x^{k+1}
    =
    x^k-\eta\nabla f(x^k),
\]
где \(\eta>0\) --- шаг метода.

\begin{algorithm}[H]
\caption{Градиентный спуск с постоянным шагом}
\label{alg:lecture10-gradient-descent}
\begin{algorithmic}[1]

\Require
функция \(f\), начальная точка \(x^0\), шаг \(\eta>0\)

\For{\(k=0,1,2,\ldots\)}

    \State
    \(g^k\gets\nabla f(x^k)\)

    \State
    \(x^{k+1}\gets x^k-\eta g^k\)

\EndFor

\Ensure
последовательность \(\{x^k\}_{k\ge0}\)

\end{algorithmic}
\end{algorithm}

Даже для простой квадратичной функции выбор шага определяет поведение
итераций.



Рассмотрим
\[
    f(x)=\frac12x^2.
\]
Тогда
\[
    \nabla f(x)=x,
    \qquad
    L=1,
\]
и градиентный спуск принимает вид
\[
    x^{k+1}
    =
    (1-\eta)x^k.
\]
Следовательно,
\[
    x^k
    =
    (1-\eta)^k x^0.
\]

Сходимость к \(x^\star=0\) происходит тогда и только тогда, когда
\[
    \abs{1-\eta}<1,
\]
то есть
\[
    0<\eta<2.
\]

Поведение метода различается в зависимости от шага:
\[
\begin{array}{c|c}
    \text{шаг} & \text{поведение}
    \\
    \hline
    0<\eta<1
    &
    \text{монотонное уменьшение } \abs{x^k}
    \\
    \eta=1
    &
    x^1=0
    \\
    1<\eta<2
    &
    \text{сходимость с чередованием знака}
    \\
    \eta=2
    &
    x^{k+1}=-x^k
    \\
    \eta>2
    &
    \abs{x^k}\to\infty
\end{array}
\]

Таким образом, слишком большой шаг приводит к колебаниям или
расходимости. Для рассматриваемой функции выбор
\[
    \eta=\frac1L=1
\]
даёт решение за одну итерацию.

\begin{remark}
Для произвольной выпуклой \(L\)-гладкой функции выбор
\(\eta=1/L\) является стандартным безопасным шагом. Полный вывод
оценки сходимости в данной лекции не проводится и будет дан при
последующем анализе градиентного метода.
\end{remark}

\section{Двойственная задача линейного программирования}

Рассмотрим прямую задачу
\[
    \begin{aligned}
        \minimize_{x\in\RR^n}
        \quad&
        c^{\mathsf T}x,
        \\
        \subjectto
        \quad&
        Ax\le b,
    \end{aligned}
    \tag{\ensuremath{P_{\mathrm{LP}}}}
\]
где
\[
    A\in\RR^{m\times n},
    \qquad
    b\in\RR^m,
    \qquad
    c\in\RR^n.
\]

Неравенство \(Ax\le b\) понимается покомпонентно. Введём двойственную
переменную
\[
    \lambda\in\RR^m_+.
\]

Функция Лагранжа равна
\[
\begin{aligned}
    \mathcal L(x,\lambda)
    &=
    c^{\mathsf T}x
    +
    \lambda^{\mathsf T}(Ax-b)
    \\
    &=
    \bigl(c+A^{\mathsf T}\lambda\bigr)^{\mathsf T}x
    -
    b^{\mathsf T}\lambda.
\end{aligned}
\]

Двойственная функция определяется как
\[
    g(\lambda)
    =
    \inf_{x\in\RR^n}
    \mathcal L(x,\lambda).
\]

Если
\[
    c+A^{\mathsf T}\lambda\ne0,
\]
то линейную по \(x\) функцию можно сделать сколь угодно малой,
выбирая направление \(x\). Поэтому
\[
    g(\lambda)=-\infty.
\]

Если же
\[
    A^{\mathsf T}\lambda+c=0,
\]
зависимость от \(x\) исчезает и
\[
    g(\lambda)=-b^{\mathsf T}\lambda.
\]

Таким образом,
\[
    g(\lambda)
    =
    \begin{cases}
        -b^{\mathsf T}\lambda,
        &
        A^{\mathsf T}\lambda+c=0,
        \\[1mm]
        -\infty,
        &
        A^{\mathsf T}\lambda+c\ne0.
    \end{cases}
\]

\begin{problem}[Двойственная задача LP]
Двойственная задача к \((P_{\mathrm{LP}})\) имеет вид
\[
    \begin{aligned}
        \maximize_{\lambda\in\RR^m}
        \quad&
        -b^{\mathsf T}\lambda,
        \\
        \subjectto
        \quad&
        A^{\mathsf T}\lambda+c=0,
        \\
        &
        \lambda\ge0.
    \end{aligned}
    \tag{\ensuremath{D_{\mathrm{LP}}}}
\]
\end{problem}

Двойственная задача также является линейной программой. В ней
\(m\) переменных и \(n\) равенств, тогда как в прямой задаче ---
\(n\) переменных и \(m\) неравенств. Поэтому условие \(m\ll n\)
уменьшает размер двойственного вектора, но само по себе ещё не
гарантирует меньшей вычислительной сложности.

\begin{remark}[Уточнение формулы из транскрипта]
Для конечности двойственной функции требуется именно равенство
\[
    A^{\mathsf T}\lambda+c=0.
\]
Здесь требуется именно равенство; знак неравенства не обеспечивает конечности двойственной функции.
\end{remark}



В распределённой задаче работа с прямой переменной может требовать
хранения и передачи вектора
\[
    x\in\RR^n.
\]
При переходе к двойственной постановке агенты работают с
\[
    \lambda\in\RR^m.
\]
Если число ограничений значительно меньше числа прямых переменных,
это уменьшает объём памяти и коммуникаций.

Однако решение двойственной задачи непосредственно даёт
\(\lambda^\star\), а не \(x^\star\). Восстановление прямого решения
требует дополнительных условий оптимальности. Для общей линейной
программы одного равенства
\[
    A^{\mathsf T}\lambda^\star+c=0
\]
недостаточно для однозначного определения \(x^\star\); необходимо
использовать прямую допустимость и условия дополняющей нежёсткости,
которые будут введены вместе с условиями Каруша--Куна--Таккера.

\begin{remark}[Сильная двойственность для LP]
Для линейной программы условие Слейтера со строгими неравенствами не
требуется. Если прямая задача допустима и её оптимальное значение
конечно, то двойственная задача также имеет оптимальное решение и
\[
    p^\star=d^\star.
\]
Аналогичное утверждение верно при перестановке прямой и двойственной
задач. При недопустимости или неограниченности одной из задач
необходимо отдельно разбирать соответствующий альтернативный случай.
\end{remark}

\section{Прямая задача SVM и стационарность}

Пусть дана строго линейно разделимая выборка, в которой представлены
оба класса:
\[
    \{(x_i,y_i)\}_{i=1}^{m},
    \qquad
    x_i\in\RR^n,
    \qquad
    y_i\in\{-1,+1\}.
\]

Жёстко-разделяющий метод опорных векторов решает задачу
\[
    \begin{aligned}
        \minimize_{w\in\RR^n,\ b\in\RR}
        \quad&
        \frac12\norm{w}^2,
        \\
        \subjectto
        \quad&
        y_i\bigl(w^{\mathsf T}x_i+b\bigr)\ge1,
        \qquad
        i=1,\ldots,m.
    \end{aligned}
    \tag{\ensuremath{P_{\mathrm{SVM}}}}
\]

Запишем ограничения в стандартной форме:
\[
    1-y_i\bigl(w^{\mathsf T}x_i+b\bigr)\le0.
\]
Введём множители
\[
    \lambda_i\ge0,
    \qquad
    i=1,\ldots,m.
\]

Функция Лагранжа имеет вид
\[
\begin{aligned}
    \mathcal L(w,b,\lambda)
    =
    \frac12\norm{w}^2
    +
    \sum_{i=1}^{m}
    \lambda_i
    \left[
        1-y_i\bigl(w^{\mathsf T}x_i+b\bigr)
    \right].
\end{aligned}
\]

Для построения двойственной функции необходимо минимизировать
\(\mathcal L\) по \(w\) и \(b\).



Минимизация по \(w\) даёт единственную точку
\[
    w(\lambda)
    =
    \sum_{i=1}^{m}\lambda_i y_i x_i,
\]
поскольку
\[
    \nabla_w\mathcal L
    =
    w-
    \sum_{i=1}^{m}\lambda_i y_i x_i.
\]
По переменной \(b\) функция Лагранжа линейна. Поэтому
\[
    \inf_{b\in\RR}\mathcal L(w,b,\lambda)>-\infty
\]
тогда и только тогда, когда
\[
    \sum_{i=1}^{m}\lambda_i y_i=0.
\]
При выполнении этого равенства зависимость от \(b\) исчезает.

Первая формула показывает, что нормаль к оптимальной разделяющей
гиперплоскости является линейной комбинацией объектов обучающей
выборки.

\section{Двойственная задача SVM и восстановление классификатора}

Подставим
\[
    w=\sum_{j=1}^{m}\lambda_jy_jx_j
\]
в функцию Лагранжа. Имеем
\[
\begin{aligned}
    \frac12\norm{w}^2
    &=
    \frac12
    \sum_{i=1}^{m}
    \sum_{j=1}^{m}
    \lambda_i\lambda_j
    y_iy_j
    x_i^{\mathsf T}x_j,
    \\
    \sum_{i=1}^{m}
    \lambda_i y_i w^{\mathsf T}x_i
    &=
    \sum_{i=1}^{m}
    \sum_{j=1}^{m}
    \lambda_i\lambda_j
    y_iy_j
    x_i^{\mathsf T}x_j.
\end{aligned}
\]
Слагаемое, содержащее \(b\), исчезает вследствие условия
\[
    \sum_{i=1}^{m}\lambda_i y_i=0.
\]

Поэтому двойственная функция равна
\[
    g(\lambda)
    =
    \sum_{i=1}^{m}\lambda_i
    -
    \frac12
    \sum_{i=1}^{m}
    \sum_{j=1}^{m}
    \lambda_i\lambda_j
    y_iy_j
    x_i^{\mathsf T}x_j.
\]

\begin{problem}[Двойственная задача SVM]
Двойственная задача жёстко-разделяющего SVM имеет вид
\[
    \begin{aligned}
        \maximize_{\lambda\in\RR^m}
        \quad&
        \sum_{i=1}^{m}\lambda_i
        -
        \frac12
        \sum_{i=1}^{m}
        \sum_{j=1}^{m}
        \lambda_i\lambda_j
        y_iy_j
        x_i^{\mathsf T}x_j,
        \\
        \subjectto
        \quad&
        \lambda_i\ge0,
        \qquad i=1,\ldots,m,
        \\
        &
        \sum_{i=1}^{m}\lambda_i y_i=0.
    \end{aligned}
    \tag{\ensuremath{D_{\mathrm{SVM}}}}
\]
\end{problem}

В матричной форме положим
\[
    K_{ij}=x_i^{\mathsf T}x_j,
    \qquad
    Y=\diag(y_1,\ldots,y_m).
\]
Тогда
\[
    g(\lambda)
    =
    \one^{\mathsf T}\lambda
    -
    \frac12
    \lambda^{\mathsf T}YKY\lambda.
\]

Матрица \(K\) является матрицей Грама и положительно
полуопределена. Следовательно, \(YKY\succeq0\), а функция
\(g\) вогнута.



Поскольку выборка строго линейно разделима, ограничения можно
выполнить строго после подходящего масштабирования \((w,b)\).
Следовательно, условие Слейтера выполнено, сильная двойственность
имеет место, а KKT-условия необходимы и достаточны.

После решения двойственной задачи нормаль гиперплоскости находится
по формуле
\[
    w^\star
    =
    \sum_{i=1}^{m}
    \lambda_i^\star y_i x_i.
\]

Значение \(b^\star\) определяется с помощью опорного вектора.
Для индекса \(i\), для которого
\[
    \lambda_i^\star>0,
\]
условие дополняющей нежёсткости даёт
\[
    y_i
    \bigl(
        w^{\star\mathsf T}x_i+b^\star
    \bigr)
    =
    1.
\]
Следовательно,
\[
    b^\star
    =
    y_i-w^{\star\mathsf T}x_i.
\]

В самой лекции это восстановление откладывается до введения условий
Каруша--Куна--Таккера.

Решающее правило имеет вид
\[
    \widehat y(x)
    =
    \sign
    \left(
        \sum_{i=1}^{m}
        \lambda_i^\star y_i
        x_i^{\mathsf T}x
        +
        b^\star
    \right).
\]

Точки с
\[
    \lambda_i^\star=0
\]
не входят в итоговую сумму. Ненулевые множители соответствуют
опорным векторам.

\section{Ядровый переход и отображения признаков}

Двойственная задача зависит от объектов \(x_i\) только через
скалярные произведения
\[
    x_i^{\mathsf T}x_j.
\]
Пусть
\[
    \phi\colon\RR^n\to\mathcal H
\]
отображает исходные признаки в другое пространство, где данные
могут стать линейно разделимыми. Тогда в двойственной задаче
возникают только величины
\[
    \inner{\phi(x_i)}{\phi(x_j)}_{\mathcal H}.
\]

\begin{definition}[Ядровая функция]
Функция
\[
    K(x,z)
    =
    \inner{\phi(x)}{\phi(z)}_{\mathcal H}
\]
называется ядром, соответствующим отображению признаков \(\phi\).
\end{definition}

В двойственной задаче достаточно выполнить замену
\[
    x_i^{\mathsf T}x_j
    \longmapsto
    K(x_i,x_j).
\]
Получаем
\[
    \begin{aligned}
        \maximize_{\lambda\in\RR^m}
        \quad&
        \sum_{i=1}^{m}\lambda_i
        -
        \frac12
        \sum_{i,j=1}^{m}
        \lambda_i\lambda_j
        y_iy_j
        K(x_i,x_j),
        \\
        \subjectto
        \quad&
        \lambda_i\ge0,
        \\
        &
        \sum_{i=1}^{m}\lambda_i y_i=0.
    \end{aligned}
\]

Классификатор принимает форму
\[
    \widehat y(x)
    =
    \sign
    \left(
        \sum_{i=1}^{m}
        \lambda_i^\star y_i
        K(x_i,x)
        +
        b^\star
    \right).
\]

Пространство признаков \(\mathcal H\) при этом не требуется строить
явно. Этот приём называется ядровым трюком.



Пусть в \(\RR^2\) классы расположены около окружностей
\[
    x_1^2+x_2^2=1
    \qquad\text{и}\qquad
    x_1^2+x_2^2=4.
\]
В исходных координатах такие классы нельзя разделить одной прямой.

Рассмотрим новый признак
\[
    z=x_1^2+x_2^2.
\]
Тогда внутренний класс располагается около значения \(z=1\), а
внешний --- около \(z=4\). В новом пространстве их можно разделить
линейным порогом
\[
    z=\beta,
    \qquad
    1<\beta<4.
\]

В исходных координатах этому порогу соответствует нелинейная граница
\[
    x_1^2+x_2^2=\beta.
\]
Следовательно, линейный классификатор в пространстве признаков задаёт
нелинейную границу в исходном пространстве.

Распространённые примеры ядер:
\[
    K_{\mathrm{poly}}(x,z)
    =
    \bigl(x^{\mathsf T}z+c\bigr)^p,
\]
\[
    K_{\mathrm{RBF}}(x,z)
    =
    \exp
    \left(
        -\frac{\norm{x-z}^2}{2\sigma^2}
    \right).
\]

\begin{remark}
Не всякая функция двух аргументов является допустимым ядром. Для
симметричной функции \(K\) необходимо и достаточно, чтобы для любой
конечной системы точек \(x_1,\ldots,x_m\) матрица Грама
\[
    \bigl[K(x_i,x_j)\bigr]_{i,j=1}^{m}
\]
была положительно полуопределённой.
\end{remark}

\section{Практическая роль двойственной формы}

Для линейной программы двойственный переход меняет размерность
оптимизируемого вектора с \(n\) на \(m\). Это может уменьшить объём
передаваемого состояния при \(m\ll n\), однако двойственная задача
одновременно содержит \(n\) равенств, поэтому вычислительное
преимущество зависит от структуры матриц и используемого алгоритма.

Для SVM основное преимущество двойственной формы состоит не только
в размерности. В ней данные входят исключительно через скалярные
произведения, что позволяет применить ядровый переход и получить
нелинейный классификатор.

При этом двойственная постановка также имеет вычислительные
ограничения. В SVM число двойственных переменных равно числу объектов
обучающей выборки, а матрица Грама имеет размер
\[
    m\times m.
\]
Поэтому при очень большом числе объектов двойственная задача может
оказаться дороже прямой.

\begin{lecturesummary}
Для градиентного спуска
\[
    x^{k+1}=x^k-\eta\nabla f(x^k)
\]
размер шага принципиально влияет на сходимость. В квадратичном
примере \(f(x)=x^2/2\) метод сходится при \(0<\eta<2\), а шаг
\(\eta=1/L\) приводит к минимуму за одну итерацию.

Для линейной программы
\[
    \minimize_x c^{\mathsf T}x
    \quad
    \subjectto
    \quad
    Ax\le b
\]
двойственная задача равна
\[
    \maximize_{\lambda\ge0}
    -b^{\mathsf T}\lambda
    \quad
    \subjectto
    \quad
    A^{\mathsf T}\lambda+c=0.
\]

Двойственная задача жёстко-разделяющего SVM имеет вид
\[
    \maximize_{\lambda\ge0}
    \left\{
        \sum_i\lambda_i
        -
        \frac12
        \sum_{i,j}
        \lambda_i\lambda_jy_iy_j
        x_i^{\mathsf T}x_j
    \right\},
    \qquad
    \sum_i\lambda_i y_i=0.
\]
Зависимость только от скалярных произведений позволяет заменить их
ядрами \(K(x_i,x_j)\) и получить нелинейную разделяющую поверхность
без явного построения пространства признаков.
\end{lecturesummary}